% GENERATED FILE. Edit canonical lessons, then run npm run generate:books.
% canonical-source-hash: fnv1a64:4ee966b5170a5b5f
% canonical-lessons: ES-C465-cuaderno, ES-C465-anotar, ES-C465-temperatura, ES-C465-lluvia, ES-R465-repaso-apunte, ES-C465-sintesis-apunte

\chapter{El cuaderno de campo}
\label{ch:es-apunte}
\hlchaptermodality{465}

\begin{chapteropening}
\textbf{By the end of this chapter:} \emph{I can keep a shared written record: read the instructions inside a weather notebook, note down a temperature and a rainfall, and understand why a zero is not the same as a blank.}

Read the instructions taped inside a shared weather notebook, then fill in a day out loud.
\end{chapteropening}

\section[el cuaderno]{el cuaderno — the notebook}

\label{lesson:ES-C465-cuaderno}



\begin{quote}
Say \textbf{el libro}, \textbf{el papel} and \textbf{el carné}. One is bound, one is loose, and the third already told you which number it counts to.
\end{quote}

\subsection*{You'll want to know: el cuaderno}
\textbf{el cuaderno} — the notebook, the exercise book.

\begin{quote}
Traiga un \textbf{cuaderno} y un bolígrafo.
\end{quote}

\emph{Un cuaderno de espiral} — a spiral notebook. \emph{El cuaderno de campo} — the field notebook, the one a person carries to write down what they see outside.

A \emph{libro} is written for you. A \emph{cuaderno} is waiting for you to write in it. That is the difference that matters, and it is why a school hands out \emph{cuadernos} and a library lends \emph{libros}.

\begin{cousinweb}
You have had this number already. \emph{El carné} was Latin \textbf{quaternum}, a set of four — four sheets of paper folded together, which is how a small notebook was made — and that lesson named \emph{el cuaderno} while it was there. This is the object the card shrank from.

\noindent
\begin{tabularx}{\linewidth}{@{}>{\raggedright\arraybackslash}X>{\raggedright\arraybackslash}X@{}}
\toprule
\textbf{Latin} & \textbf{Spanish} \\
\midrule
quaternum & el carné, el cuaderno \\
quattuor & cuatro, cuarto \\
\bottomrule
\end{tabularx}

So \emph{quaternum} counts \textbf{sheets gathered into a fold}, not pages produced by one. Four leaves folded together was the cheapest book anybody could make, and the thing was named after how many went in.

English kept the count plainer still: a \textbf{quire} is a gathering of folded sheets, and it is the same \emph{quaternum} worn down through French.

The notebook in a bag and the number four are one word, separated by the thing rarely thought about when buying one — how many sheets went into the fold.
\end{cousinweb}

\begin{grammarlens}[title={Grammar Lens: the ua that does not move}]
Spanish has taught you to watch stressed vowels, because those are the ones that break. This family shows the other half of that rule.

\noindent
\begin{tabularx}{\linewidth}{@{}>{\raggedright\arraybackslash}X>{\raggedright\arraybackslash}X@{}}
\toprule
\textbf{word} & \textbf{stress falls on} \\
\midrule
\textbf{cua}-tro & the \emph{ua} \\
cua-\textbf{der}-no & the \emph{-der-} \\
\bottomrule
\end{tabularx}

The \emph{ua} sits still in both. It is not a broken vowel waiting to happen — it came straight from the Latin \emph{qua-} and stayed. The lesson is that the stress test tells you where to \emph{look}, and looking sometimes tells you nothing is happening. Breaking needs two things, not one: the vowel must be stressed \textbf{and} must have been a short Latin \textbf{e} or \textbf{o} (or the \emph{ae} that fell in with short \emph{e}, which is why \emph{quaerere} gave \emph{quiero}). In \emph{cuatro} the \emph{ua} is stressed, so it passes the first test and fails the second. In \emph{cuaderno} it fails both.
\end{grammarlens}

\subsection*{Your turn}
Say these aloud:

\begin{itemize}
\raggedright
  \item bring a notebook
  \item the difference between a libro and a cuaderno
  \item which number is hiding in cuaderno
\end{itemize}

\begin{tcolorbox}[breakable,colback=teal!4,colframe=teal!35!black,title={Before you move on}]
The notebook? (\textbf{El cuaderno}.) Which word did you meet its number under? (\emph{\textbf{El carné}} — the same \emph{quaternum}, four sheets in one fold.)
\end{tcolorbox}
\section[anotar]{anotar — to write it down}

\label{lesson:ES-C465-anotar}



\begin{quote}
Say \textbf{apuntar} and \textbf{la nota}. You already have a verb for writing something down. A second one is about to arrive, and the two lean in different directions.
\end{quote}

\subsection*{You'll want to know: anotar}
\textbf{anotar} — to note down, to record in writing.

\begin{quote}
\textbf{Anote} la temperatura cada mañana en el cuaderno.
\end{quote}

\emph{Anotar los datos} — to record the figures. \emph{Anotar un gol} — to score a goal, because the scoreboard is a record too.

\begin{grammarlens}[title={Grammar Lens: anotar is not apuntar}]
Both put words on paper, and the difference is what the writing is \emph{for}.

\noindent
\begin{tabularx}{\linewidth}{@{}>{\raggedright\arraybackslash}X>{\raggedright\arraybackslash}X@{}}
\toprule
\textbf{what you are doing} & \textbf{the verb} \\
\midrule
jotting so you do not forget & \textbf{apuntar} \\
entering it into a record & \textbf{anotar} \\
\bottomrule
\end{tabularx}

\emph{Apunta mi número} is a scrap of paper in a pocket. \emph{Anote su número en el registro} is a line in a book that will be read later by somebody else.

The etymologies lean the same way, which is rare enough to be worth noticing. \emph{Apuntar} is \textbf{punctum}, the dot a point leaves behind — you had it from \emph{pungere}, to pierce. \emph{Anotar} is \textbf{nota}, a mark made so that somebody could read it. Both name marks; the difference is who each mark is for.

So a doctor \emph{apunta} while you talk and \emph{anota} in your file afterwards.
\end{grammarlens}

\begin{cousinweb}
Latin \textbf{adnotare}: \textbf{ad-}, towards, plus \textbf{nota}, the mark you have from \emph{la nota}. The writing is all in the noun.

Latin \emph{nota} was a sign that stood for something — a letter, a brand, a mark in the margin. English took it both ways: a \textbf{note} you write, and to \textbf{notice}, which marks a thing without writing anything.

You were given \textbf{notar} alongside \emph{la nota} — to notice, to catch the mark. That is Spanish having taken the second road. \emph{Anotar} is the first road, with the prefix aiming the mark at something.

\noindent
\begin{tabularx}{\linewidth}{@{}>{\raggedright\arraybackslash}X>{\raggedright\arraybackslash}X@{}}
\toprule
\textbf{Spanish} & \textbf{what it does} \\
\midrule
notar & to notice it \\
anotar & to write it down \\
\bottomrule
\end{tabularx}

One root, one prefix, and the prefix is the whole difference between seeing something and keeping it.

\textbf{Anotado} is the participle, and \emph{las anotaciones} are the notes themselves — the \emph{-ción} ending doing the job it has done since you first met it.
\end{cousinweb}

\subsection*{Your turn}
Say these aloud:

\begin{itemize}
\raggedright
  \item note the temperature down in the notebook
  \item which verb you use for a scrap of paper
  \item what notar means without the a-
\end{itemize}

\begin{tcolorbox}[breakable,colback=teal!4,colframe=teal!35!black,title={Before you move on}]
To write it into a record? (\textbf{Anotar}.) And to notice, without writing? (\textbf{Notar}.) Which Latin word is inside both? (\emph{\textbf{Nota}}.)
\end{tcolorbox}
\section[la temperatura]{la temperatura — the temperature}

\label{lesson:ES-C465-temperatura}



\begin{quote}
Say \textbf{el calor} and \textbf{el frío}. You have the two ends. Here is the word for the scale that runs between them.
\end{quote}

\subsection*{You'll want to know: la temperatura}
\textbf{la temperatura} — the temperature.

\begin{quote}
Anote la \textbf{temperatura} a las ocho de la mañana y a las ocho de la tarde.
\end{quote}

\emph{Tomar la temperatura} — to take somebody's temperature. In Spain, \emph{tener temperatura} is to be running a fever with no number said out loud; \emph{tener fiebre} is the safer phrase everywhere.

Note what \emph{calor} and \emph{frío} cannot do: they tell you which way, never how much. \emph{Hace calor} is a judgement. \emph{La temperatura} is a measurement, and the whole reason a notebook wants one is that a measurement can be compared with yesterday's.

\begin{cousinweb}
Latin \textbf{temperare} did not mean heat at all. It meant \textbf{to mix in the right proportion} — to hold two things in balance so that neither wins.

Follow that one idea and a very odd family falls into line:

\noindent
\begin{tabularx}{\linewidth}{@{}>{\raggedright\arraybackslash}X>{\raggedright\arraybackslash}X@{}}
\toprule
\textbf{English} & \textbf{what is being balanced} \\
\midrule
to \textbf{temper} steel & heat against hardness \\
a \textbf{temperate} climate & summer against winter \\
\bottomrule
\end{tabularx}

A person's \textbf{temper} is the mix of their humours, and losing it is the mix coming apart. \textbf{Tempera} paint is pigment held in the right proportion of egg. And a \textbf{temperature} is the point a mixture has reached on the scale between hot and cold.

Spanish keeps the verb as \textbf{templar}, to temper or to tune an instrument, and \emph{templado} is the word for weather that is neither one thing nor the other.
\end{cousinweb}

\begin{grammarlens}[title={Grammar Lens: -ura has two jobs}]
\emph{La factura} gave you \textbf{-ura} as the thing an action produces, with a verb hiding inside the noun. \emph{Temperatura} is built the same way and on the same instruction: the verb inside it is \emph{temperare}, and what the action produces is a reading.

The ending has a second family, though, and it does not take verbs at all:

\noindent
\begin{tabularx}{\linewidth}{@{}>{\raggedright\arraybackslash}X>{\raggedright\arraybackslash}X@{}}
\toprule
\textbf{base} & \textbf{noun} \\
\midrule
alto & la alt\textbf{ura} \\
ancho & la anch\textbf{ura} \\
\bottomrule
\end{tabularx}

Those sit on adjectives, and they name how much of the quality there is — \emph{altura} is how tall, \emph{anchura} is how wide, and \emph{la verdura} is the one you already own. So when you meet a new \textbf{-ura} noun, look for a verb first, as \emph{factura} told you to, and look for an adjective second.
\end{grammarlens}

\subsection*{Your turn}
Say these aloud:

\begin{itemize}
\raggedright
  \item note down the temperature
  \item what temperare meant before it meant heat
  \item how tall, using the -ura noun
\end{itemize}

\begin{tcolorbox}[breakable,colback=teal!4,colframe=teal!35!black,title={Before you move on}]
The temperature? (\textbf{La temperatura}.) What did \emph{temperare} mean? (\textbf{To mix in the right proportion}.) What does \emph{-ura} build? (\textbf{The thing an action produces} — \emph{factura}, \emph{temperatura} — and on an adjective, \textbf{how much of a quality} — \emph{altura}.)
\end{tcolorbox}
\section[la lluvia]{la lluvia — the rain}

\label{lesson:ES-C465-lluvia}



\begin{quote}
Say \textbf{llueve} and \textbf{el llanto}. Both begin with the sound this word begins with, and both have already told you why.
\end{quote}

\subsection*{You'll want to know: la lluvia}
\textbf{la lluvia} — the rain.

\begin{quote}
\textbf{Lluvia} en los últimos tres días: 12 mm.
\end{quote}

\emph{Lluvia fuerte} — heavy rain. \emph{Un día de lluvia} — a rainy day. \emph{La lluvia ácida} — acid rain.

\begin{grammarlens}[title={Grammar Lens: the verb and the noun do different work}]
You already have \emph{llueve}, and it refuses a subject: there is nobody raining. \emph{La lluvia} is what you reach for the moment rain has to be counted, described or put in a column.

\noindent
\begin{tabularx}{\linewidth}{@{}>{\raggedright\arraybackslash}X>{\raggedright\arraybackslash}X@{}}
\toprule
\textbf{what you are saying} & \textbf{the word} \\
\midrule
it is raining now & \textbf{llueve} \\
how much rain there was & \textbf{la lluvia} \\
\bottomrule
\end{tabularx}

\emph{Llueve mucho} is a report from the window. \emph{Hubo mucha lluvia} is a line in a notebook. A verb with no subject cannot be measured; a noun can, which is why the rainfall column is headed \emph{lluvia} and never \emph{llueve}.
\end{grammarlens}

\begin{cousinweb}
This word's Latin has been passed under your nose several times already, and it is worth saying so.

The rule came to you under \emph{llamar}, back at the beginning: Latin \emph{cl-}, \emph{fl-} and \emph{pl-} arrive as Spanish \textbf{ll-} in words that came down by mouth. \emph{Llueve} put it on rain, from \emph{pluere}. And under \emph{el llanto} it was run out over a list — where \emph{la lluvia} was already sitting, as the example beside \emph{lleno} and \emph{llorar}.

\noindent
\begin{tabularx}{\linewidth}{@{}>{\raggedright\arraybackslash}X>{\raggedright\arraybackslash}X@{}}
\toprule
\textbf{Latin} & \textbf{Spanish} \\
\midrule
pluvia & la lluvia \\
plenus & lleno \\
\bottomrule
\end{tabularx}

So nothing here is new except the use. Latin \textbf{pluvia} is the rain itself, \emph{pluere} is the raining, and both came down by mouth with the cluster shifted.

The shift is not a law, and you have seen it fail. Words lifted later out of writing keep their \emph{pl-} — which is why \emph{planta} still has one, and why Spanish owns \textbf{la lluvia} and \textbf{pluvial} at the same time, the same word at two speeds. A \emph{pluviómetro} measures exactly what this chapter's notebook is for, and it kept the cluster because it was built by scholars, not worn down by mouths.

That is what the earlier lessons were for. A rule you were given at the start turns up later holding a word you now need, and the word costs you nothing because the rule already paid for it.
\end{cousinweb}

\subsection*{Your turn}
Say these aloud:

\begin{itemize}
\raggedright
  \item heavy rain, and a rainy day
  \item which one goes in a column, llueve or la lluvia
  \item what Latin pl- turns into
\end{itemize}

\begin{tcolorbox}[breakable,colback=teal!4,colframe=teal!35!black,title={Before you move on}]
The rain? (\textbf{La lluvia}.) Its Latin form? (\emph{\textbf{Pluvia}}.) And why does it start with \emph{ll}? (\textbf{Because Latin \emph{pl-} became Spanish \emph{ll-}}.)
\end{tcolorbox}
\section[(el cuaderno, anotar, la temperatura …]{(el cuaderno, anotar, la temperatura, la lluvia) — from cold}

\label{lesson:ES-R465-repaso-apunte}



\begin{quote}
Four words, no prompts. The book you write in, the verb for writing into it, the number you write, and the weather you write it about.
\end{quote}

\subsection*{You'll want to know: the four, cold}
\noindent
\begin{tabularx}{\linewidth}{@{}>{\raggedright\arraybackslash}X>{\raggedright\arraybackslash}X@{}}
\toprule
\textbf{English} & \textbf{Spanish} \\
\midrule
the notebook & \textbf{el cuaderno} \\
to note down & \textbf{anotar} \\
the temperature & \textbf{la temperatura} \\
the rain & \textbf{la lluvia} \\
\bottomrule
\end{tabularx}

\begin{grammarlens}[title={Grammar Lens: an impression is not a record}]
This chapter is about the moment a thing stops being an impression and becomes a record.

\noindent
\begin{tabularx}{\linewidth}{@{}>{\raggedright\arraybackslash}X>{\raggedright\arraybackslash}X@{}}
\toprule
\textbf{impression} & \textbf{record} \\
\midrule
\emph{hace calor} & \emph{la temperatura: 24} \\
\emph{llueve} & \emph{lluvia: 12 mm} \\
\bottomrule
\end{tabularx}

\emph{Hace calor} and \emph{llueve} are both true and neither can be compared with yesterday. The nouns can. That is why the chapter's verb is \textbf{anotar} rather than \textbf{apuntar}: a note you jot is for you, and a note you enter is for whoever reads the book next.
\end{grammarlens}

\begin{cousinweb}
Two of the four hide something that is not about writing or weather at all.

\noindent
\begin{tabularx}{\linewidth}{@{}>{\raggedright\arraybackslash}X>{\raggedright\arraybackslash}X@{}}
\toprule
\textbf{word} & \textbf{what is underneath} \\
\midrule
el cuaderno & \textbf{four} — \emph{quaternum}, four sheets in one fold \\
la temperatura & \textbf{mixing} — \emph{temperare}, to hold in proportion \\
\bottomrule
\end{tabularx}

The other two say what they are. \emph{Anotar} is \emph{nota}, a mark, with a prefix aiming it. And \emph{la lluvia} is \emph{pluvia}, which had been passed under your nose several times — under \emph{llueve}, \emph{el llanto}, \emph{lleno} and \emph{planta} — before you needed it.
\end{cousinweb}

\subsection*{Your turn}
Say these aloud:

\begin{itemize}
\raggedright
  \item note the temperature down in the notebook
  \item the rain, and then it is raining
  \item which of the four counts to four
\end{itemize}

\begin{tcolorbox}[breakable,colback=teal!4,colframe=teal!35!black,title={Before you move on}]
The notebook? (\textbf{El cuaderno}.) To enter it in the record? (\textbf{Anotar}.) The temperature? (\textbf{La temperatura}.) The rain? (\textbf{La lluvia}.)
\end{tcolorbox}
\section[Practice]{(el cuaderno de campo) — read it, then fill in a day}

\label{lesson:ES-C465-sintesis-apunte}



\begin{quote}
Say \textbf{el cuaderno} and \textbf{anotar}. Somebody else will read what you write in it, which is the whole reason the instructions exist.
\end{quote}

\begin{tcolorbox}[breakable,skin=enhanced,colback=orange!6,colframe=orange!45!black,boxrule=0.5pt,arc=1mm,left=6pt,right=6pt,top=4pt,bottom=4pt,fonttitle=\bfseries,title={Read this}]
Taped inside the front cover of a shared notebook in a community garden:

\begin{quote}
\textbf{CUADERNO DE CAMPO — CÓMO USARLO}
\end{quote}

\begin{quote}
\textbf{Anote} la temperatura dos veces al día: a las ocho de la mañana y a las ocho de la tarde.
\end{quote}

\begin{quote}
Si ha llovido, \textbf{anote} también la lluvia en milímetros. Si no ha llovido, escriba un cero. No deje la casilla en blanco.
\end{quote}

\begin{quote}
Una línea por día. Si olvida un día, déjelo vacío y siga: es peor \textbf{anotar} de memoria que no \textbf{anotar} nada.
\end{quote}
\end{tcolorbox}

\begin{grammarlens}[title={Grammar Lens: a zero is not a blank}]
Four instructions, and the last one is the interesting one.

\noindent
\begin{tabularx}{\linewidth}{@{}>{\raggedright\arraybackslash}X>{\raggedright\arraybackslash}X@{}}
\toprule
\textbf{the instruction} & \textbf{what it wants} \\
\midrule
\emph{anote la temperatura} & twice a day, at fixed hours \\
\emph{escriba un cero} & a zero, not an empty box \\
\bottomrule
\end{tabularx}

The notice draws a line between \textbf{un cero} and \textbf{la casilla en blanco}. A zero says it did not rain. A blank says nobody looked. Those are different facts, and a record that cannot tell them apart is worth less than one that can — which is why the last line would rather you left a gap than filled it in from memory.
\end{grammarlens}

\begin{grammarlens}[title={Grammar Lens: anote, not anota}]
Every instruction here is \textbf{usted}: \emph{anote}, \emph{escriba}, \emph{deje}, \emph{siga}. A notice usually takes the polite form, because it does not know who will read it — though plenty of Spanish notices use a bare infinitive instead, as in \emph{no fumar}.

\noindent
\begin{tabularx}{\linewidth}{@{}>{\raggedright\arraybackslash}X>{\raggedright\arraybackslash}X@{}}
\toprule
\textbf{to one person you know} & \textbf{on a notice} \\
\midrule
\emph{anota} & \textbf{anote} \\
\bottomrule
\end{tabularx}

You have had the mechanism since \emph{hable usted}: an \emph{usted} command \textbf{is} the subjunctive, because \emph{usted} is a third-person title. Not a borrowing — the same form doing its own job. (The \emph{tú} negative is subjunctive too: \emph{no anotes}.) It is why \emph{reduzca} on the road sign looked the way it did.
\end{grammarlens}

\subsection*{Your turn}
Yesterday it rained in the morning and the afternoon was warm. Fill in the line out loud: say the date, the two temperatures, and the rain.

Now say what you would write if it had not rained at all — and say why that is not the same as leaving it blank.

\begin{tcolorbox}[breakable,colback=teal!4,colframe=teal!35!black,title={Before you move on}]
The field notebook? (\textbf{El cuaderno de campo}.) The polite command to note something? (\textbf{Anote}.) What goes in the box on a dry day? (\textbf{Un cero} — never a blank.)
\end{tcolorbox}
